% Paragraph goals and sentence sub-goals; intentionally not polished prose.
\section*{Current paper outline after the September 4 meeting}
\label{sec:sept-outline}
\noindent\textbf{Scope.} Study practical memory selection and observable memory
auditing, building on existing discovery methods. Retain simulation as background.
Do not claim identification of hidden intentions or a true real-world latent SCM.

\noindent\textbf{Execution status, September 15.} The faithful comparison is not
complete. Three online mechanism checks and two Travel development cases passed;
both authorized metered relays currently report insufficient balance. Full-history
LoCoMo validation is incomplete and the 100 formal cases have not started.
P10--P13 describe experiment goals, not completed new results.

\paragraph{P1: Motivate memory quality.}
S1: Agents depend on earlier records. S2: Full histories cost context, while
retrieval and summaries can lose dependencies or propagate errors.
S3: A correct visible answer can coexist with a risky retrieved record.

\paragraph{P2: State the empirical question.}
S1: Connect writes, reads and query targets. S2: Ask whether an estimated temporal
structure helps select context and trace observable risks. S3: Test utility under
a shared model and the original benchmark actor without assuming real-world identifiability.

\paragraph{P3: Match contributions to evidence.}
S1: Present the traceable memory pipeline. S2: Compare task utility and total memory
cost against recent native baselines. S3: Provide a concrete risk case.
S4: Keep negative mitigation results and identification limits visible.

\paragraph{P4: Define variables across regions or regimes.}
S1: Each environment has its own schema. S2: Travel shares seven slot types, while
traveler/day/record counts vary. S3: Regime-specific edges do not change column
semantics. S4: The transcript's region/regime wording remains uncertain.

\paragraph{P5: Define samples and allowed information.}
S1: A traveler round is time; an episode is a trial. S2: The frozen graph uses query
constraint indicators and within-trial lag windows. S3: Online memory contains the
public base and actor writebacks. S4: Held-out gold is reserved for evaluation.

\paragraph{P6: Separate evaluation questions.}
S1: P2 measures downstream utility and cost. S2: P3-A traces candidate historical
drivers. S3: P3-B tests intervention and retained task utility.
S4: None directly measures an unobserved psychological intention.

\paragraph{P7: Use existing discovery as a component.}
S1: Specify observed inputs and lag-graph outputs. S2: Cite assumptions and existing
simulation results. S3: Interpret real-data outputs as candidate dependencies.

\paragraph{P8: Explain the memory operation.}
S1: Parse query target cells. S2: Combine estimated type edges with explicit instance
references. S3: Retrieve source values through write provenance.
S4: Expose the full public base as memory; the original actor emits a full plan.
No local inheritance decoder supplies answer cells.

\paragraph{P9: Explain auditing and labels.}
S1: Trace observable write/read/action paths. S2: Separate oracle-tagged MINJA
illustrations from label-free AgentPoison calibration. S3: Freeze driver and gate
rules. S4: Distinguish attribution from effective mitigation.

\paragraph{P10: Introduce recognizable baselines.}
S1: Describe MemoryArena's multi-round task. S2: Include full context, BM25, dense
retrieval and rolling autoregressive summaries. S3: Add Mem0 OSS, the authors'
AgenticMemory paper-repository robust agent and the complete LightMem core pipeline.
S4: Validate full-history LoCoMo task interfaces separately; do not equate task
conformance with published-score or managed-product replication.

\paragraph{P11: Fix comparison and accounting.}
S1: The current development uses aiaaa deepseek-v4-flash-0731 and the original actor, tools and full-plan
output. S2: Develop on 101; freeze 111--120 only after implementation gates pass.
These historical holdout IDs are reused.
S3: Report official PS/SPS/SR and episode-level values. S4: Count memory-writing,
updating and actor tokens and time; estimates are not relay invoices.
S5: Use task-evaluation retrieval sizes 200/60/10 for Mem0/LightMem/A-Mem,
preserve A-Mem's native 1000-token output cap, and disable memory reasoning with
the verified relay parameter. Do not pool results from different relay models.

\paragraph{P12: Interpret P2 results after completion.}
S1: Fill the new paired table from artifacts. S2: Explain quality/cost tradeoffs even
when a baseline wins. S3: Keep the historical ten-episode result separate.
S4: The old noG comparison also changes representation; the new noGcompact holds
format fixed and disables selection. The query-only control preserves explicit
references but removes learned type edges. S5: Reused episodes and one API
realization cannot establish a stable ranking or general noninferiority.

\paragraph{P13: Show one concrete memory slide.}
S1: Show a real query and required cells. S2: Distinguish estimated type edges,
parsed references and selected records. S3: Compare actual contexts and answers
with rolling summaries and a native baseline. S4: Do not invent a correction if
all methods succeed or the proposed method fails.

\paragraph{P14: Show dormant risk without mind-reading.}
S1: MINJA round 15 writes record fixed\_1\_3 containing an answer-shift rule.
S2: Round 16 retrieves it but answers correctly. S3: Round 19 has no attack note,
retrieves it alongside other records and returns H instead of D.
S4: Explain the estimated gated edge and oracle labels; this is not a matched
single-record counterfactual.

\paragraph{P15: Preserve the full audit outcome.}
S1: Report cross-seed attribution evidence. S2: Explain label-free selection of
AgentPoison record 9251 and cluster expansion to 9252. S3: Show removal without
recovered task correctness. S4: Retain both failed online mitigation protocols.

\paragraph{P16: Position related work.}
S1: Compare extraction/update systems. S2: Compare RAG and summaries.
S3: Build on causal discovery rather than repeating its simulation story.
S4: Distinguish source auditing from robust online defense.

\paragraph{P17: State actionable limitations.}
S1: Disclose explicit query references and the visible public base; remove the
local domain decoder from the new comparison.
S2: Discuss variable choice, confounding and finite trials. S3: Restrict cross-task
claims. S4: Disclose reused holdouts, one API realization and relay differences.
S5: Conclude only supported findings; leave generalization and mitigation open.

\noindent\textbf{Figures and tables.} Prioritize a concrete memory-selection slide
and dormant-risk slide; then a setting table, paired baseline task/cost table and
separate attribution/mitigation table. Existing simulations and full traces belong
in the appendix. The Chinese companion gives the same paragraph/sentence goals.

\noindent\textbf{Historical exploratory result, September 14.} The ten-arm paired table on three
reused holdout episodes is complete. Nine arms use the original frozen v5 outputs;
the interrupted rolling-summary arm uses the first complete technical rerun under
identical source code and settings, with all attempts retained for cost accounting.
Ours and same-format noGcompact tie on every episode's PS/SPS/SR. Ours uses 49.63\%
fewer input tokens, but its usage-based total cost is 4.59\% higher because of
additional output. This supports an observed input reduction, not extra accuracy,
monetary savings or causal necessity. The Travel slide remains a matched-history
read replay; the MINJA risk slide uses stored observable trace events.

\noindent\textbf{Remediation, September 15 (in progress).} The historical v5
comparison does not meet the complete-implementation requirement and is excluded
from the new main table. Core mechanism tests passed for Mem0 OSS, paper A-Mem
and LightMem. Native-task and original-actor validation are in progress; new
paired results remain pending. Preserve actor token-limit failures unchanged and
report all failed/development usage.
